% from ega1-10-fr.tex
\oldpage[I]{182}
\begin{env}[10.1.5]
\label{I.10.1.5-fr}
En tant que faisceau d'anneaux \emph{sans topologie}, le faisceau
structural $\mathscr{O}_{\mathfrak{X}}$ de
$\operatorname{Spf}(A)$ admet pour tout $x\in\mathfrak{X}$ une fibre qui,
en vertu de (\hyperref[I.10.1.4-fr]{10.1.4}), s'identifie à la limite
inductive
$\varinjlim_{f\notin\mathfrak{j}_x}A_{\{f\}}$. Par suite
(\hyperref[0.7.6.17-fr]{0, 7.6.17} et
\hyperref[0.7.6.18-fr]{7.6.18})~:
\end{env}

\begin{proposition}[10.1.6]
\label{I.10.1.6-fr}
Pour tout $x\in\mathfrak{X}=\operatorname{Spf}(A)$, la fibre
$\mathscr{O}_x$ est un anneau local dont le corps résiduel est isomorphe à
$k(x)=A_x/\mathfrak{j}_xA_x$. Si en outre $A$ est adique et noethérien,
$\mathscr{O}_x$ est un anneau noethérien.
\end{proposition}

Comme $k(x)$ n'est pas réduit à 0, on conclut de ce résultat que le
\emph{support} du faisceau d'anneaux $\mathscr{O}_{\mathfrak{X}}$ est
\emph{égal à $\mathfrak{X}$}.

\subsection{Morphismes de schémas formels affines.}
\label{subsection:I.10.2-fr}

\begin{env}[10.2.1]
\label{I.10.2.1-fr}
Soient $A$, $B$ deux anneaux admissibles, et soit
$\varphi:B\to A$ un homomorphisme \emph{continu}. L'application continue
$ {}^a\varphi:\operatorname{Spec}(A)\to\operatorname{Spec}(B)$
(\hyperref[I.1.2.1-fr]{1.2.1}) applique alors
$\mathfrak{X}=\operatorname{Spf}(A)$ dans
$\mathfrak{Y}=\operatorname{Spf}(B)$, car l'image réciproque par
$\varphi$ d'un idéal premier ouvert de $A$ est un idéal premier ouvert de
$B$. D'autre part, pour tout $g\in B$, $\varphi$ définit un homomorphisme
continu
\[
  \Gamma(\mathfrak{D}(g),\mathscr{O}_{\mathfrak{Y}})
  \longrightarrow
  \Gamma(\mathfrak{D}(\varphi(g)),\mathscr{O}_{\mathfrak{X}})
\]
en vertu de (\hyperref[I.10.1.4-fr]{10.1.4}),
(\hyperref[I.10.1.3-fr]{10.1.3}) et
(\hyperref[0.7.6.7-fr]{0, 7.6.7})~; comme ces homomorphismes satisfont aux
conditions de compatibilité pour les restrictions correspondant au
passage de $g$ à un multiple de $g$, et que
$\mathfrak{D}(\varphi(g))={}^a\varphi^{-1}(\mathfrak{D}(g))$, ils
définissent un homomorphisme \emph{continu} de faisceaux d'anneaux
topologiques
$\mathscr{O}_{\mathfrak{Y}}\to
{}^a\varphi_*(\mathscr{O}_{\mathfrak{X}})$
(\hyperref[0.3.2.5-fr]{0, 3.2.5}), que nous noterons encore
$\widetilde{\varphi}$~; on a ainsi obtenu un morphisme
$\Phi=({}^a\varphi,\widetilde{\varphi})$ d'espaces topologiquement annelés
$\mathfrak{X}\to\mathfrak{Y}$. On notera qu'en tant qu'homomorphisme de
faisceaux d'anneaux sans topologie, $\widetilde{\varphi}$ définit un
homomorphisme
$\widetilde{\varphi}_x^\sharp:
\mathscr{O}_{{}^a\varphi(x)}\to\mathscr{O}_x$ sur les fibres, pour tout
$x\in\mathfrak{X}$.
\end{env}

\begin{proposition}[10.2.2]
\label{I.10.2.2-fr}
Soient $A$, $B$ deux anneaux topologiques admissibles, et soient
$\mathfrak{X}=\operatorname{Spf}(A)$,
$\mathfrak{Y}=\operatorname{Spf}(B)$. Pour qu'un morphisme
$u=(\psi,\theta):\mathfrak{X}\to\mathfrak{Y}$ d'espaces topologiquement
annelés soit de la forme $({}^a\varphi,\widetilde{\varphi})$, où
$\varphi$ est un homomorphisme continu d'anneaux $B\to A$, il faut et il
suffit que pour tout $x\in\mathfrak{X}$,
$\theta_x^\sharp$ soit un homomorphisme local
$\mathscr{O}_{\psi(x)}\to\mathscr{O}_x$.
\end{proposition}

\begin{proof}
La condition est nécessaire~: en effet, soit
$\mathfrak{p}=\mathfrak{j}_x\in\operatorname{Spf}(A)$, et soit
$\mathfrak{q}=\varphi^{-1}(\mathfrak{j}_x)$~; si
$g\notin\mathfrak{q}$, on a donc $\varphi(g)\notin\mathfrak{p}$, et il est
immédiat que l'homomorphisme
$B_{\{g\}}\to A_{\{\varphi(g)\}}$ déduit de $\varphi$
(\hyperref[0.7.6.7-fr]{0, 7.6.7}) transforme
$\mathfrak{q}_{\{g\}}$ en une partie de
$\mathfrak{p}_{\{\varphi(g)\}}$~; passant à la limite inductive, on voit
donc (compte tenu de (\hyperref[I.10.1.5-fr]{10.1.5}) et de
(\hyperref[0.7.6.17-fr]{0, 7.6.17})) que
$\widetilde{\varphi}_x^\sharp$ est un homomorphisme local.

Inversement, soit $(\psi,\theta)$ un morphisme vérifiant la condition de
l'énoncé~; en vertu de (\hyperref[I.10.1.3-fr]{10.1.3}), $\theta$ définit
un homomorphisme continu d'anneaux
\[
  \varphi=\Gamma(\theta):
  B=\Gamma(\mathfrak{Y},\mathscr{O}_{\mathfrak{Y}})
  \longrightarrow
  \Gamma(\mathfrak{X},\mathscr{O}_{\mathfrak{X}})=A.
\]
En vertu de l'hypothèse sur $\theta$, pour que la section $\varphi(g)$ de
$\mathscr{O}_{\mathfrak{X}}$ au-dessus de $\mathfrak{X}$ ait un germe
inversible au point $x$, il faut et il suffit que $g$ ait un germe
inversible au point $\psi(x)$. Mais en vertu de
(\hyperref[0.7.6.17-fr]{0, 7.6.17}), les sections de
$\mathscr{O}_{\mathfrak{X}}$ (resp. $\mathscr{O}_{\mathfrak{Y}}$)
au-dessus de $\mathfrak{X}$ (resp. $\mathfrak{Y}$) dont le germe n'est pas
inversible au point $x$ (resp. $\psi(x)$) sont exactement les éléments de
$\mathfrak{j}_x$
